\documentclass[aspectratio=169]{beamer}

\usetheme{Madrid}
\usecolortheme{seahorse}
\usepackage[utf8]{inputenc}
\usepackage[T1]{fontenc}
\usepackage[polish]{babel}
\usepackage{tikz}
\usepackage{amsmath}
\usepackage{animate}
\usetikzlibrary{arrows.meta, shapes.geometric, positioning}

\title[Traffic Circle MPI]{Symulacja ronda metodą Monte Carlo z MPI}
\author{Systemy Równoległe}
\date{\today}

% ==================================================================
\begin{document}
% ==================================================================

\begin{frame}
  \titlepage
\end{frame}

% ==================================================================
\begin{frame}{Problem}
  \begin{columns}
    \begin{column}{0.52\textwidth}
      Rondo z $R$ wjazdami i $R$ wyjazdami.\\[0.4em]
      Pytania (podręcznik, rozdz. 10.5.6):
      \begin{enumerate}
        \item[(a)] $P(\text{czekanie})$ — prawdopodobieństwo, że auto trafi do kolejki
        \item[(b)] $\overline{q}$ — średnia długość kolejki przy wjeździe
      \end{enumerate}
      \bigskip
      Metoda: \textbf{dyskretna symulacja Monte Carlo}\\
      z równoległym uruchamianiem replik przez MPI.
    \end{column}
    \begin{column}{0.44\textwidth}
      \begin{tikzpicture}[scale=0.9]
        \draw[thick, gray!50, fill=gray!10] (0,0) circle (1.8cm);
        \draw[thick, white, fill=white] (0,0) circle (0.7cm);
        % TikZ uzywa stopni bezposrednio
        \foreach \angle/\lbl/\col in {90/N/blue, 0/E/orange, 270/S/red, 180/W/teal}{
          % droga: linia od zewnatrz do krawedzi kola
          \draw[\col, very thick] (\angle:1.85) -- (\angle:2.5);
          % strzalka wjazdu (do srodka)
          \draw[\col, -latex, thick] (\angle:2.5) -- (\angle:2.0);
          % etykieta
          \node[\col!80!black, font=\bfseries] at (\angle:3.0) {\lbl};
        }
        \draw[-latex, gray, thick] (55:1.3) arc[start angle=55, end angle=15, radius=1.3cm];
        \node[font=\tiny, gray] at (0,0) {$4R$ slotów};
      \end{tikzpicture}

      \centering\small Rondo dla $R=4$ (16 slotów)
    \end{column}
  \end{columns}
\end{frame}

% ==================================================================
\begin{frame}{Model symulacji}
  Rondo podzielone na $4R$ slotów. Wjazd i wyjazd drogi $i$ — ten sam slot $i \cdot 4$.\\[0.6em]

  \textbf{Każda iteracja — 3 fazy:}
  \begin{enumerate}
    \item \textbf{Przyjazdy} — dla każdego wjazdu $i$: odlicz czas; jeśli minął → nowe auto (czas z $\mathrm{Exp}(f_i)$)
    \item \textbf{Ruch na rondzie} — wszystkie auta przesuwają się o 1 slot; auto na swoim slocie docelowym opuszcza rondo
    \item \textbf{Wjazd na rondo} — slot wolny: wjeżdża auto z kolejki lub nowe; slot zajęty: nowe auto czeka w kolejce
  \end{enumerate}

  \bigskip
  \textbf{Rozgrzewka:} pierwsze 10\% iteracji odrzucane — ustala się stan ustalony.\\[0.3em]
  \textbf{Parametry} ($R=4$): $f = [3, 3, 4, 2]$, macierz $D[4][4]$ z rys. 10.21.
\end{frame}

% ==================================================================
\begin{frame}{Paralelizacja — replikacja Monte Carlo}
  \begin{columns}
    \begin{column}{0.55\textwidth}
      Każdy proces MPI uruchamia \textbf{niezależną kopię} całej symulacji.

      \bigskip
      \begin{itemize}
        \item Iteracje dzielone równomiernie: $\ell_k = \lfloor N/P \rfloor$
        \item Niezależny generator LCG na każdy rank\\
              \texttt{seed = (rank+1) * C \^{} time(NULL)}
        \item Wyniki agregowane: \texttt{MPI\_Reduce(SUM)}
        \item Rank 0 oblicza i wyświetla średnie ważone
      \end{itemize}

      \bigskip
      Złożoność komunikacji: $O(R)$ liczb — minimalna.
    \end{column}
    \begin{column}{0.4\textwidth}
      \begin{tikzpicture}[node distance=0.45cm, font=\small]
        \node[draw, rounded corners, fill=blue!15, minimum width=3.2cm, minimum height=0.55cm] (r0) {rank 0 \ \ $\ell_0$ iteracji};
        \node[draw, rounded corners, fill=blue!15, minimum width=3.2cm, minimum height=0.55cm, below=of r0] (r1) {rank 1 \ \ $\ell_1$ iteracji};
        \node[below=0.05cm of r1, font=\small] {$\vdots$};
        \node[draw, rounded corners, fill=blue!15, minimum width=3.2cm, minimum height=0.55cm, below=0.65cm of r1] (rn) {rank $P{-}1$ \ $\ell_{P-1}$ iteracji};
        \node[draw, rounded corners, fill=orange!30, minimum width=3.2cm, minimum height=0.6cm, below=0.5cm of rn] (red) {MPI\_Reduce(SUM)};
        \node[draw, rounded corners, fill=green!20, minimum width=3.2cm, minimum height=0.55cm, below=0.4cm of red] (out) {rank 0: wyniki / $N$};
        \draw[-latex] (r0.south) -- ++(0,-0.12) -| (red.north);
        \draw[-latex] (r1.south) -- ++(0,-0.12) -| (red.north);
        \draw[-latex] (rn.south) -- (red.north);
        \draw[-latex] (red.south) -- (out.north);
      \end{tikzpicture}
    \end{column}
  \end{columns}
\end{frame}

% ==================================================================
\begin{frame}{$P$ (procesy MPI) i $R$ (wjazdy) — niezależne parametry}
  \begin{columns}
    \begin{column}{0.50\textwidth}
      \textbf{Geometria ronda:}
      \[
        \text{sloty} = 4 \times R
      \]
      Liczba slotów jest \textbf{zawsze} 4-krotnością liczby wjazdów.\\
      Wynika to z odstępu między drogami w modelu.

      \bigskip
      \textbf{$P$ i $R$ są od siebie niezależne:}
      \begin{itemize}
        \item $R$ określa \textbf{model} — ile dróg ma rondo
        \item $P$ określa \textbf{zasoby} — ile węzłów klastra użyć
        \item Każda kombinacja $(P, R)$ jest poprawna
      \end{itemize}

      \bigskip
      Każdy z $P$ procesów symuluje \textbf{to samo rondo} ($R$ wjazdów)\\
      przez $N/P$ iteracji — więcej węzłów = więcej próbek = \\
      mniejsza wariancja wyników.
    \end{column}
    \begin{column}{0.46\textwidth}
      \begin{center}
      \begin{tabular}{cccc}
        \hline
        $R$ & sloty ($4R$) & przykładowe $P$ \\
        \hline
        4  & 16 & 1, 2, 4, 8, 16, \ldots \\
        8  & 32 & 1, 2, 4, 8, 16, \ldots \\
        16 & 64 & 1, 2, 4, \ldots, 64 \\
        \hline
      \end{tabular}
      \end{center}

      \bigskip
      \begin{tikzpicture}[scale=0.82]
        \draw[-latex] (0,0) -- (5.0,0) node[right, font=\small] {$R$};
        \draw[-latex] (0,0) -- (0,4.6) node[above, font=\small] {$P$};
        \foreach \x in {1,2,3,4}{
          \draw[gray!25] (\x,0) -- (\x,4.3);
          \draw[gray!25] (0,\x) -- (4.8,\x);
          \draw (\x,.05) -- (\x,-.05) node[below,font=\scriptsize] {\x};
          \draw (.05,\x) -- (-.05,\x) node[left, font=\scriptsize] {\x};
        }
        % P = k*R lines
        \draw[blue!60,  thick] (0,0) -- (4.3,4.3) node[right,font=\scriptsize,blue!60]  {$P{=}R$};
        \draw[teal!70,  thick] (0,0) -- (2.15,4.3) node[above,font=\scriptsize,teal!70] {$P{=}2R$};
        \draw[red!60,   thick] (0,0) -- (1.075,4.3) node[above,font=\scriptsize,red!60] {$P{=}4R$};
        % area — dowolne P
        \fill[yellow!20, opacity=0.5] (0,0) -- (4.8,0) -- (4.8,4.3) -- (0,4.3) -- cycle;
        \node[font=\scriptsize, gray] at (3.0,2.0) {dowolne $P$};
      \end{tikzpicture}
    \end{column}
  \end{columns}
\end{frame}

% ==================================================================
\begin{frame}{Uruchomienie}
  \begin{block}{Kompilacja}
    \texttt{mpicc -O2 -o traffic\_circle\_mpi traffic\_circle\_mpi.c -lm}
  \end{block}

  \begin{block}{Uruchomienie}
    \texttt{mpirun -np \textit{P} ./traffic\_circle\_mpi \textit{iteracje} \textit{R} [\texttt{--anim}]}
  \end{block}

  \bigskip
  \begin{center}
    \begin{tabular}{lll}
      \hline
      Parametr & Domyślnie & Opis \\
      \hline
      \textit{P}         & —        & liczba procesów MPI \\
      \textit{iteracje}  & 500\,000 & łączna liczba kroków symulacji \\
      \textit{R}         & 4        & liczba wjazdów/wyjazdów (max 16) \\
      \texttt{--anim}    & wył.     & zapis stanu ronda do CSV (tylko $R=4$) \\
      \hline
    \end{tabular}
  \end{center}

  \bigskip
  \textbf{Przykład} (16 węzłów klastra, 1M iteracji, 4 drogi):
  \begin{center}
    \texttt{mpirun -np 16 --hostfile nodes ./traffic\_circle\_mpi 1000000 4}
  \end{center}
\end{frame}

% ==================================================================
% Aby ten slajd działał, przed kompilacją wyeksportuj klatki z GIFa:
%   mkdir frames
%   convert -coalesce traffic_animation.gif frames/frame%04d.png
% Animacja działa w Adobe Acrobat; w innych przeglądarkach PDF widoczna
% jest tylko pierwsza klatka.
% ==================================================================
\begin{frame}{Wizualizacja działania ronda ($R=4$)}
  \begin{columns}
    \begin{column}{0.62\textwidth}
      \centering
      \animategraphics[loop, autoplay, width=\textwidth]{15}{frames/frame}{0000}{0599}
    \end{column}
    \begin{column}{0.35\textwidth}
      \small
      \textbf{Kolory} = docelowy wyjazd auta:\\[0.4em]
      {\color{red!70!black}$\blacksquare$} N \quad
      {\color{blue!70!black}$\blacksquare$} W \quad
      {\color{green!60!black}$\blacksquare$} S \quad
      {\color{orange!80!black}$\blacksquare$} E \\[0.8em]
      \textbf{Słupki} przy wjazdach\\= długość kolejki\\[0.8em]
      Każda klatka = 1 iteracja\\
      symulacji (stan ustalony).\\[0.8em]
      Generowanie:\\[0.2em]
      \texttt{\scriptsize mpirun -np 16 ./traffic}\\
      \texttt{\scriptsize \ \ 500000 4 --anim}\\
      \texttt{\scriptsize python3 animate\_traffic.py}
    \end{column}
  \end{columns}
\end{frame}

% ==================================================================
\begin{frame}{}
  \centering
  \Huge Dziękuję za uwagę
\end{frame}

\end{document}
